\section{Teorema de Taylor}

Como hemos observado anteriormente, si $f$ es diferenciable en $c$, entonces $f$ es aproximadamente una función lineal en las proximidades de $c$. Esto es, la ecuación
\[
f(x)=f(c)+f'(c)(x-c)
\]
es aproximadamente correcta cuando $(x-c)$ es pequeño. Pero, ¿qué pasa cuando conocemos la segunda derivada o la tercera derivada, o incluso si $f$ tiene derivada n-ésima?

Pues bien, el teorema de Taylor nos dice que, en general, $f$ puede aproximarse por medio de un polinomio de grado $n-1$, si $f$ posee derivadas hasta de orden $n$. Además, el teorema de Taylor proporciona una expresión útil para calcular el error cometido en esta aproximación.

El teorema de Taylor se obtiene como consecuencia de un resultado más general; veamos entonces el siguiente teorema.

\begin{teorema}
    Sean $f$ y $g$ dos funciones que poseen derivadas n-ésimas $f^{(n)}$ y $g^{(n)}$ finitas en un intervalo abierto $(a,b)$ y derivadas $(n-1)$ continuas en el intervalo cerrado $[a,b]$. Supongamos que $c\in[a,b]$. Entonces, para todo $x\in[a,b]$, $x\neq c$, existe un punto $x_{1}$ interior al intervalo que une a $x$ con $c$ tal que
    \[
    \left[ f(x)-\sum_{k=0}^{n-1}\frac{f^{(k)}(c)}{k!}(x-c)^{k} \right]g^{(n)}(x_{1})
    =
    f^{(n)}(x_{1})\left[ g(x)-\sum_{k=0}^{n-1}\frac{g^{(k)}(c)}{k!}(x-c)^{k} \right].
    \]
\end{teorema}

\begin{prueba}
    Supongamos que $c\in [a,b]$ y tomemos $x\geq c$ con $x$ fijo (el caso $x<c$ es análogo). Definamos las funciones auxiliares $F(t)$ y $G(t)$ para cada $t\in[c,x]$ como
    \[
    F(t)= f(t)+\sum_{k=1}^{n-1}\frac{f^{(k)}(t)}{k!}(x-t)^{k},
    \qquad
    G(t)= g(t)+\sum_{k=1}^{n-1}\frac{g^{(k)}(t)}{k!}(x-t)^{k}.
    \]
    \begin{itemize}
        \item Cuando $t=x$, la suma se anula porque $(x-x)=0$; entonces obtenemos $F(x)=f(x)$ y $G(x)=g(x)$.
        \item Cuando $t=c$, tenemos
        \[
        G(c)= g(c)+\sum_{k=1}^{n-1}\frac{g^{(k)}(c)}{k!}(x-c)^{k}
        = g(c)+g'(c)(x-c)+\frac{g^{(2)}(c)}{2!}(x-c)^{2}+\dots+\frac{g^{(n-1)}(c)}{(n-1)!}(x-c)^{n-1}
        = \sum_{k=0}^{n-1}\frac{g^{(k)}(c)}{k!}(x-c)^{k},
        \]
        y, de manera análoga,
        \[
        F(c)= \sum_{k=0}^{n-1}\frac{f^{(k)}(c)}{k!}(x-c)^{k},
        \]
        que son precisamente los polinomios de Taylor de grado $n-1$.
    \end{itemize}
    Aplicamos el Teorema del Valor Medio Generalizado, que nos indica que existe $x_{1}\in(c,x)$ tal que
    \[
    F'(x_{1})\left[ G(x)-G(c) \right]=G'(x_{1})\left[ F(x)-F(c) \right]. \tag{$*$}
    \]
    Sustituimos $F(x)=f(x)$ y $G(x)=g(x)$:
    \[
    F'(x_{1})\left[ g(x)-G(c) \right]=G'(x_{1})\left[ f(x)-F(c) \right].
    \]
    Calculamos ahora $F'(t)$ y $G'(t)$. Derivando $F(t)$ (para $G(t)$ es análogo), para un término con $k$ fijo y recordando que $x$ es constante:
    \[
    \frac{d}{dt}\left[ \frac{f^{(k)}(t)}{k!}(x-t)^{k} \right]
    =\frac{f^{(k+1)}(t)}{k!}(x-t)^{k}-\frac{f^{(k)}(t)}{(k-1)!}(x-t)^{k-1}.
    \]
    Sumamos para $k=1$ hasta $n-1$:
    \[
    F'(t)=f'(t)+\sum_{k=1}^{n-1}\frac{f^{(k+1)}(t)}{k!}(x-t)^{k}-\sum_{k=1}^{n-1}\frac{f^{(k)}(t)}{(k-1)!}(x-t)^{k-1}.
    \]
    Hacemos un cambio de índice en la segunda sumatoria: tomamos $j=k-1$, con $j=0$ hasta $n-2$:
    \[
    \sum_{k=1}^{n-1}\frac{f^{(k)}(t)}{(k-1)!}(x-t)^{k-1}
    =\sum_{j=0}^{n-2}\frac{f^{(j+1)}(t)}{j!}(x-t)^{j},
    \]
    y así
    \[
    F'(t)=f'(t)+\sum_{k=1}^{n-1}\frac{f^{(k+1)}(t)}{k!}(x-t)^{k}-\sum_{j=0}^{n-2}\frac{f^{(j+1)}(t)}{j!}(x-t)^{j}.
    \]
    Observemos que el término $f'(t)$ corresponde a $j=0$ en la última suma. Por tanto, al restar, todos los términos desde $j=0$ hasta $n-2$ se cancelan, subsistiendo únicamente el último término de la primera suma, cuando $k=n-1$; obteniendo así
    \[
    F'(t)=\frac{f^{(n)}(t)}{(n-1)!}(x-t)^{n-1}.
    \]
    Análogamente,
    \[
    G'(t)=\frac{g^{(n)}(t)}{(n-1)!}(x-t)^{n-1}.
    \]
    Recordando que $t=x_{1}$, y sustituyendo $F'(x_{1})$ y $G'(x_{1})$ en $(*)$, tenemos
    \[
    \frac{f^{(n)}(x_{1})}{(n-1)!}(x-x_{1})^{n-1}\left[ g(x)-G(c) \right]
    =
    \frac{g^{(n)}(x_{1})}{(n-1)!}(x-x_{1})^{n-1}\left[ f(x)-F(c) \right],
    \]
    de donde, cancelando el factor $\frac{(x-x_{1})^{n-1}}{(n-1)!}\neq 0$,
    \[
    f^{(n)}(x_{1})\left[ g(x)-G(c) \right]=g^{(n)}(x_{1})\left[ f(x)-F(c) \right].
    \]
    Reemplazando $F(c)$ y $G(c)$ por los polinomios de Taylor de grado $n-1$, obtenemos
    \[
    \left[ f(x)-\sum_{k=0}^{n-1}\frac{f^{(k)}(c)}{k!}(x-c)^{k} \right]g^{(n)}(x_{1})
    =
    f^{(n)}(x_{1})\left[ g(x)-\sum_{k=0}^{n-1}\frac{g^{(k)}(c)}{k!}(x-c)^{k} \right],
    \]
    que es lo que queríamos probar.
\end{prueba}

\begin{teorema}
     Sea $f$ una función que admite derivada n-ésima $f^{(n)}$ finita en todo el intervalo abierto $(a,b)$ y supongamos que $f^{(n-1)}$ es continua en el intervalo cerrado $[a,b]$. Supongamos que $c\in[a,b]$. Entonces, para todo $x\in[a,b]$, $x\neq c$, existe un punto $x_{1}$, interior al intervalo que une $x$ con $c$, tal que
     \[
     f(x)=f(c)+\sum_{k=1}^{n-1}\frac{f^{(k)}(c)}{k!}(x-c)^{k}+\frac{f^{(n)}(x_{1})}{n!}(x-c)^{n}.
     \]
\end{teorema}

\begin{prueba}
    Este teorema es un caso particular del teorema anterior si consideramos $g(x)=(x-c)^{n}$. Dadas $f$ y $g$ funciones con derivadas n-ésimas finitas en $(a,b)$ y derivadas $(n-1)$-ésimas continuas en $[a,b]$, para todo $x\in[a,b]$, $x\neq c$, existe $x_{1}$ en el intervalo abierto que une $x$ y $c$ tal que
    \begin{equation}\label{taylor}
    \left[ f(x)-\sum_{k=0}^{n-1}\frac{f^{(k)}(c)}{k!}(x-c)^{k} \right]g^{(n)}(x_{1})=f^{(n)}(x_{1})\left[ g(x)-\sum_{k=0}^{n-1}\frac{g^{(k)}(c)}{k!}(x-c)^{k} \right].
    \end{equation}
    De $g(x)=(x-c)^{n}$ se cumple que $g(c)=0$, $g'(c)=0$, \dots, $g^{(k)}(c)=0$ para todo $k<n$. Por lo tanto,
    \[
    \sum_{k=0}^{n-1}\frac{g^{(k)}(c)}{k!}(x-c)^{k}=0,
    \]
    de modo que el lado derecho de \eqref{taylor} se reduce a
    \begin{equation}\label{tay1}
    f^{(n)}(x_{1})(x-c)^{n}.
    \end{equation}
    Ahora calculamos la derivada n-ésima de $g$. Para $g(x)=(x-c)^{n}$:
    \[
    g'(x)=n(x-c)^{n-1},
    \qquad
    g^{(2)}(x)=n(n-1)(x-c)^{n-2},
    \qquad
    \dots
    \]
    y, en general,
    \[
    g^{(k)}(x)=n(n-1)(n-2)\dots(n-k+1)(x-c)^{n-k}.
    \]
    Cuando $k=n$, $(x-c)^{0}=1$, quedando así el producto
    \[
    n(n-1)(n-2)\dots 2\cdot 1 = n!,
    \]
    por lo tanto, para cualquier punto $x_{1}$,
    \begin{equation}\label{tay2}
    g^{(n)}(x_{1})=n!.
    \end{equation}
    Sustituyendo \eqref{tay1} y \eqref{tay2} en \eqref{taylor} se tiene
    \[
    \left[ f(x)-\sum_{k=0}^{n-1}\frac{f^{(k)}(c)}{k!}(x-c)^{k} \right]n! = f^{(n)}(x_{1})(x-c)^{n},
    \]
    luego
    \[
    f(x)-\sum_{k=0}^{n-1}\frac{f^{(k)}(c)}{k!}(x-c)^{k} =\frac{f^{(n)}(x_{1})}{n!}(x-c)^{n}.
    \]
    Así,
    \[
    f(x)=\sum_{k=0}^{n-1}\frac{f^{(k)}(c)}{k!}(x-c)^{k}+\frac{f^{(n)}(x_{1})}{n!}(x-c)^{n},
    \]
    que es lo que queríamos probar.
\end{prueba}
